%% SCRAP: scratch/testing-quality/KANBAN-WORKFLOW-VALIDATION
%% SOURCE: docs/working/scratch/testing-quality/KANBAN-WORKFLOW-VALIDATION.adoc
%% STATUS: OBSOLETE
%% FITS: none
%% EDITORIAL: lifted — prose rewritten to press voice

\section{Kanban Workflow Validation Procedure (2025-11-02)}

This document recorded the formal validation procedure for the GitHub Projects V2
Kanban workflow automation implemented as Phase 3 of CAPA-032. It is preserved as a
historical test-procedure artefact; the automation it describes is no longer the
current state of the project.

\subsection{Workflow Under Test}

File: \texttt{.github/workflows/capa-kanban-sync.yml}. The workflow automated six
status transitions for CAPA issues: CREATE, IMPLEMENT, VALIDATE, APPROVE, RELEASE,
and CLOSED. Transitions were triggered by pull-request events, workflow completions,
and comment patterns.

\subsection{Validation Summary}

The procedure required all five prerequisite checks to pass (GitHub CLI version,
authentication, repository access, project existence, and workflow deployment) before
test execution. The full lifecycle test drove a CAPA issue through all six states by
creating a PR, waiting for the test workflow, posting QA and PM approval comments, and
merging the PR. Success required all 13 checklist items to pass with no manual
interventions beyond the two approval comments. Expected total elapsed time was
approximately 20 minutes.

Artifact integrity was verified by SHA-256 checksums:
\begin{itemize}
  \item \texttt{KANBAN-SETUP-CHECKLIST.sh}: \texttt{ad9e585538aa320521e6b4af449ead6a83c9c2471e1ae2d4d592792d9bb0deaa}
  \item \texttt{KANBAN-QUICK-TEST.sh}: \texttt{ef13d408853a1bd221b4f31a8968a729f14ba1f9d02d5bff82cf733a39070a20}
\end{itemize}

%% TODO(bob): confirm whether this validation was completed and whether the workflow
%% remains active or has been superseded.
